\begin{figure}[t]
  \centering
  \includegraphics[width=\textwidth]{figures/HippoRAG_hook_fig_version_4.pdf}
  \caption{
  \textbf{Knowledge Integration \& RAG.} Tasks that require knowledge integration are particularly challenging for current RAG systems. In the above example, we want to find a \textit{Stanford} professor that does \textit{Alzheimer's} research from a pool of passages describing potentially thousands \textit{Stanford} professors and \textit{Alzheimer's} researchers.
  Since current methods encode passages in isolation, they would struggle to identify \textit{Prof. Thomas} unless a passage mentions both characteristics at once.
  In contrast, most people familiar with this professor would remember him quickly due to our brain's associative memory capabilities, thought to be driven by the index structure depicted in the C-shaped hippocampus above (in \textcolor{hippoblue}{blue}).
  Inspired by this mechanism, \textbf{HippoRAG} allows LLMs to build and leverage a similar graph of associations to tackle knowledge integration tasks.
  }
  \label{fig:hook}
  \vspace{-15pt}
\end{figure}
